%! Author = WelJo
%! Date = 4/30/2026

%    The abstract should briefly summarize the contents of the paper in
%    150--250 words~\cite{ref_book1}.

\begin{abstract}
    Critical infrastructures (CIs), such as hospitals, fire stations, and utility networks, are vital for urban functionality yet remain highly vulnerable to extreme environmental hazards.
    While existing research frameworks typically evaluate either transport-related road accessibility or utility failure cascades in isolation, they fail to capture how cascading outages dynamically degrade spatial emergency service delivery.
    In this paper, we propose a dynamic framework linking the Robustness of Accessibility (RoA) index with time-dependent modeling of cascading infrastructure failures.
    We extend the RoA metric across a continuous disaster horizon and introduce a four-state Finite State Machine (FSM), spanning operational, depleting, dead, and rebooting states, to govern facility operationality based on power and water grid dependencies, backup buffer capacities, and recovery hysteresis.
    Interdependent utility assets along subterranean corridors are topologically coupled to critical Points of Interest (POIs) using Network-Constrained Nearest-Neighbor (NCNN) graph routing.
    We implement a decoupled, two-phase geospatial simulation pipeline and evaluate a 336-hour HQ100 flood scenario in Trier, Germany, across four complexity tiers for hospitals and fire stations.
    The empirical results confirm that transport-only models substantially overestimate urban disaster resilience, yielding a 6.38\% higher cumulative accessibility index than the compound failure model.
    Power outages emerge as the primary systemic cascading bottleneck, whereas backup buffers grant critical operational bridging during flood onset.
    Furthermore, facility reboot latencies induce a pronounced 48-hour post-flood recovery lag, and topological coupling exposes severe vulnerabilities in physically unflooded peripheral districts due to upstream utility inundation.
    Ultimately, our findings provide actionable insights for municipal disaster management and climate-resilient infrastructure planning.

    %    \keywords{First keyword \and Second keyword \and Another keyword.}

    \keywords{Critical Infrastructure \and Cascading Failures \and Robustness of Accessibility \and Finite State Machine \and Urban Resilience \and Flood Hazard \and Network-Constrained Routing.}
\end{abstract}

\newpage